\documentclass[10pt]{article}

\usepackage{graphicx}
\usepackage[utf8]{inputenc}
\usepackage[spanish]{babel}
\usepackage{minted}
\usepackage{amsmath}
\usepackage[a4paper,top=2cm,bottom=2cm,left=3cm,right=3cm, marginparwidth=1.75cm]{geometry}
\usepackage{amsfonts}
\DeclareMathAlphabet{\pazocal}{OMS}{zplm}{m}{n}
\newcommand{\Lb}{\pazocal{L}}


\usemintedstyle{default}

\title{Práctica 2: Autómatas Finitos}
\author{Miguel Ángel Dorado Maldonado}

\begin{document}
\maketitle
Considera el lenguaje sobre el alfabeto \{a,b\} que solamente contiene la cadena a: L = \{a\}.
\section{Autómata finito en JFLAP.}

\subsection{Construye un AFD que reconozca el lenguaje y rechace todas las cadenas que no pertenezcan al lenguaje.}
\begin{center}
    \centering
    \includegraphics[width=0.5\linewidth]{1.1afd.png}
\end{center}
\subsection{Verifica el autómata construido con 6 cadenas de ejemplo.}
Probaremos con las siguientes cadenas:
a, ab, b, aa, ba, abababa.

\begin{center}
    \centering
    \includegraphics[width=0.5\linewidth]{1.2test.png}
\end{center}

\section{Autómata finito en Python}
\subsection{Abre el módulo automata y crea en Python un autómata equivalente al construido en la actividad anterior mediante finite\_automaton.}
Creamos el autómata en el archivo json siguiendo la estructura definida. En este caso el autómata creado lo nombraré "solo cadena a afd".
\newpage

\begin{minted}[frame=single,
               framesep=3mm,
               xleftmargin=21pt,
               tabsize=4]{js}
    	{
		"name": "solo cadena a afd",
		"representation": {
			"K": ["q0", "q1", "q2"],
			"A": ["a", "b"],
			"s": "q0",
			"F": ["q1"],
			"t": [["q0", "a", "q1"],
				  ["q0", "b", "q2"],
				  ["q1", "a", "q2"],
				  ["q1", "b", "q2"],
				  ["q2", "a", "q2"],
				  ["q2", "b", "q2"]]
		}
	}
\end{minted}


\begin{minted}
from python.automata import finite_automaton
from python.util import load_representation

# ej1
automaton = load_representation(
    "libs/python/automata/finiteautomata", "solo cadena a afd"
)
finite_automaton(automaton, "a")
finite_automaton(automaton, "ab")
finite_automaton(automaton, "b")
\end{minted}
Salida:

\begin{center}
    \centering
    \includegraphics[width=0.9\linewidth]{salidaafd.png}
\end{center}
\section{Repite los pasos anteriores pero con un AFND y compara su comportamiento con el AFD}
\subsection{Apartado 1.}
\begin{center}
    \centering
    \includegraphics[width=0.5\linewidth]{3afnd.png}
\end{center}

\begin{center}
    \centering
    \includegraphics[width=0.5\linewidth]{image.png}
\end{center}
\subsection{Apartado 2.}

\begin{minted}[frame=single,
               framesep=3mm,
               xleftmargin=21pt,
               tabsize=4]{js}
{
		"name": "solo cadena a afnd",
		"representation": {
			"K": ["q0", "q1" ],
			"A": ["a", "b"],
			"s": "q0",
			"F": ["q1"],
			"t": [["q0", "a", "q1"]]
		}
	}
\end{minted}


\begin{minted}
from python.automata import finite_automaton
from python.util import load_representation

# ej1
automaton = load_representation(
    "libs/python/automata/finiteautomata", "solo cadena a afd"
)
finite_automaton(automaton, "a")
finite_automaton(automaton, "ab")
finite_automaton(automaton, "b")
\end{minted}
Salida:
\begin{center}
    \centering
    \includegraphics[width=0.8\linewidth]{afndexample.png}
\end{center}
\end{document}

